%==============================================================================
% FICHAMENTO CRÍTICO — Use of the Deep Learning Approach to Measure Alveolar
% Bone Level
% Conforme NBR 6023 (referências) e NBR 10520 (citações)
%==============================================================================
\section{Fichamento: \citeonline{lee2022alveolar}}

% --- 1. IDENTIFICAÇÃO ---
\subsection{Identificação}
\begin{tabular}{ll}
\textbf{Título:} & Use of the deep learning approach to measure alveolar bone
level \\
\textbf{Autor(es):} & Lee, Chun-Teh; Kabir, Tanjida; Nelson, Jiman; Sheng, Sally;
Meng, Hsiu-Wan; Van Dyke, Thomas E.; Walji, Muhammad F.; Jiang, Xiaoqian;
Shams, Shayan \\
\textbf{Periódico:} & Journal of Clinical Periodontology \\
\textbf{Ano:} & 2022 \\
\textbf{Volume:} & 49 \\
\textbf{Número:} & 3 \\
\textbf{Páginas:} & 260--269 \\
\textbf{DOI:} & \url{https://doi.org/10.1111/jcpe.13574} \\
\textbf{Qualis:} & A1 (Odontologia) \\
\textbf{Tipo:} & Artigo original \\
\end{tabular}

% --- 2. OBJETIVO ---
\subsection{Objetivo}
Desenvolver e validar um modelo de deep learning baseado em redes neurais
convolucionais para medir automaticamente o nível ósseo alveolar radiográfico
em radiografias periapicais, atribuir estágios de perda óssea radiográfica
(RBL) conforme a classificação de periodontite de 2018 e comparar o desempenho
com examinadores humanos independentes.

% --- 3. METODOLOGIA ---
\subsection{Metodologia}
\textbf{Desenho do estudo:} Estudo transversal de desenvolvimento e validação de
modelo CAD baseado em DL. \\
\textbf{Amostra:} Radiografias periapicais de 46 casos (644 imagens adicionais
para validação externa independente), além do conjunto principal de treino
(70\%), validação (10\%) e teste (20\%). \\
\textbf{Métricas principais:} Dice Similarity Coefficient (DSC) para
segmentação; AUROC, sensibilidade, especificidade e acurácia para estadiamento;
comparação de porcentagem de RBL por DL versus examinadores (teste t pareado). \\
\textbf{Validação:} Três redes de segmentação integradas (osso alveolar, dente,
junção cemento-esmalte) seguidas de análise de imagem; validação em dataset
adicional independente (644 imagens de 46 casos distintos).

% --- 4. RESULTADOS PRINCIPAIS ---
\subsection{Resultados Principais}
\begin{itemize}
  \item DSC médio para segmentação foi superior a 0,91 para todas as estruturas.
  \item AUROC para estadiamento RBL: estágio I = 0,89; estágio II = 0,90;
        estágio III = 0,90; ausência de perda = 0,98.
  \item Acurácia do diagnóstico do caso (periodontite) foi de 0,85, com
        acurácia por extensão (0,96), estágio (0,87) e grau (0,96).
  \item Não houve diferença significativa entre DL e examinadores na medição
        de RBL (p = 0,65), mas o DL foi 15 vezes mais rápido (9 min vs. 137 min).
\end{itemize}

% --- 5. LIMITAÇÕES ---
\subsection{Limitações}
\begin{itemize}
  \item A amostra para validação do diagnóstico (46 casos) é limitada, com
        distribuição desigual dos estágios periodontais.
  \item O modelo foi desenvolvido apenas para radiografias periapicais, não
        tendo sido testado em radiografias panorâmicas ou CBCT.
  \item A classificação utilizou apenas o componente radiográfico, sem integrar
        dados clínicos (sondagem, sangramento, mobilidade).
\end{itemize}

% --- 6. CONTRIBUIÇÃO PARA O LIVRO ---
\subsection{Contribuição para o Livro}
Este artigo fundamenta o Capítulo 14, que trata de deep learning para
periodontia. A abordagem de três redes de segmentação integradas (osso, dente,
JCE) para medir o nível ósseo alveolar é um exemplo didático de arquitetura
modular baseada em CNNs. A comparação direta com examinadores humanos (p = 0,65)
e a economia de tempo (9 min vs. 137 min) são evidências usadas no livro para
discutir o valor clínico da automação em periodontia. A classificação conforme
a classificação de 2018 alinha-se à prática clínica contemporânea.

% --- 7. ANÁLISE CRÍTICA ---
\subsection{Análise Crítica}
\begin{itemize}
  \item \textbf{Validade interna:} Alta -- Desenho controlado com divisão
        treino/validação/teste e dataset adicional independente para evitar
        data snooping.
  \item \textbf{Validade externa:} Média -- Amostra de um único centro
        (UTHealth, EUA), mas com validação em dataset adicional de 46 casos
        distintos dos originais.
  \item \textbf{Reprodutibilidade:} Alta -- Arquitetura descrita em detalhes;
        três segmentadores CNN com parâmetros e pipeline de pós-processamento
        claramente especificados.
  \item \textbf{Relevância clínica:} Alta -- Redução de 137 para 9 minutos por
        exame completo; potencial para triagem populacional e auditoria de
        qualidade diagnóstica.
  \item \textbf{Conflito de interesses:} Não -- Os autores declararam ausência de
        conflitos.
  \item \textbf{Viés identificado:} Viés de verificação -- o ground truth foi
        estabelecido por examinadores que também serviram como comparador;
        apesar do cuidado, não houve padrão-ouro histológico.
\end{itemize}

% --- 8. CITAÇÕES RELEVANTES ---
\subsection{Citações Relevantes}
\begin{citacao}
``The average Dice Similarity Coefficient (DSC) for segmentation was over 0.91.
There was no significant difference in the RBL percentage measurements determined
by DL and examiners (p = .65).'' (p. 265).
\end{citacao}

% --- 9. MAPEAMENTO SDD ---
\subsection{Mapeamento SDD}
\textbf{Problema:} Medição manual do nível ósseo alveolar em radiografias
periapicais é subjetiva, demorada e com variabilidade interexaminador. \\
\textbf{Entrada:} Radiografias periapicais digitalizadas (intraorais) com
posicionador padronizado. \\
\textbf{Saída:} Porcentagem de RBL por sítio (mesial/distal), estágio RBL
(0--III) e diagnóstico periodontal provisório (extensão, estágio, grau). \\
\textbf{Critérios:} DSC $\geq$ 0,90; AUROC $\geq$ 0,85 por estágio; diferença
não significativa versus examinadores (p $>$ 0,05); tempo de processamento
$\leq$ 10 min por série completa. \\
\textbf{Confiança:} 0,82 -- Estudo robusto com validação em dataset independente,
mas limitado por amostra moderada e ausência de padrão-ouro histológico.

% --- 10. REFERÊNCIA ABNT ---
\subsection{Referência ABNT}
\begin{verbatim}
LEE, Chun-Teh et al. Use of the deep learning approach to measure alveolar
bone level. Journal of Clinical Periodontology, Copenhagen, v. 49, n. 3,
p. 260-269, 2022. DOI: 10.1111/jcpe.13574.
\end{verbatim}
